\section{Proof Sketch}\label{sec:proof-sketch}

The argument has five parts.  First, base conditioning gives
$\|\Lambda_M\|_{\mathrm{op}}\leq\kappa^2$ in each mode and therefore
$\|Q\|_{\mathrm{op}}\leq\kappa^6$.  Tensor-product linearity and the
product-one gauge then yield the exact coefficient identity
\[
Q(T-\widehat T_t)=D_r+E_\rho-C_t,
\qquad C_t\in\mathcal S_t.
\]

Second, orthogonal-projector geometry gives the one-step comparison
\[
\operatorname{dist}_F(D_r,\mathcal S_{t+1})
\geq
\operatorname{dist}_F(D_r,\mathcal S_t)
-\|P_{t+1}-P_t\|_{\mathrm{op}}\|D_r\|_F.
\]
The first two certificate clauses telescope this inequality over every finite
horizon and leave the all-time reserve
$(\delta-L_P)\|D_r\|_F$.

Third, distance to a fixed subspace is $1$-Lipschitz.  The smoothing residual
therefore costs at most $\zeta\|D_r\|_F$.  The exact coordinate identity,
$\|Q\|_{\mathrm{op}}\leq\kappa^6$, and
$\|T\|_F\leq C_T\|D_r\|_F$ give, for every $t$,
\[
\|T-\widehat T_t\|_F
\geq
\frac{\delta-L_P-\zeta}{\kappa^6C_T}\|T\|_F.
\]

Fourth, the third certificate clause makes the represented tensors Cauchy:
the norm of every tail difference is bounded by the corresponding tail of the
absolutely summable increment series.  Hence $\widehat T_t$ converges, and
continuity of squared Frobenius distance gives a finite objective limit.

Finally, the residual floor is squared only after checking its nonnegative
margin, and the resulting all-time scalar lower bound passes to the existing
finite limit.  This produces the event inclusion without conditioning on the
event's probability.

For the separate deterministic zero-smoothing baseline, take tall
column-orthonormal bases and set every perturbation vector to zero.  Then
$QT=D_r$, $E_\rho=0$, $\|Q\|_{\mathrm{op}}=1$, and
$\|T\|_F=\|D_r\|_F$.  Thus the same-target identity and the operator
comparison give
\[
\|T-\widehat T_t\|_F
\geq\|Q(T-\widehat T_t)\|_F
=\|D_r-C_t\|_F
\geq(\delta-L_P)\|T\|_F.
\]
Squaring and passing to the finite-variation limit yields the stronger
baseline factor $(\delta-L_P)^2$.  This algebraic baseline lies outside the
positive-Gaussian-smoothing probability statement above.
